\documentclass[11pt]{article}
\usepackage[T1]{fontenc}
\usepackage[utf8]{inputenc}
\usepackage{lmodern}
\usepackage[margin=1in]{geometry}
\usepackage{amsmath,amssymb,amsthm}
\usepackage{booktabs}
\usepackage{microtype}
\usepackage[hidelinks]{hyperref}
\hypersetup{
  pdftitle={The mixed point spectrum of a matrix contraction on its closed unit ball},
  pdfauthor={Alec Kriebel},
  pdfsubject={A self-contained completion of the mixed closed-ball Koopman point-spectrum converse}
}
\newtheorem{theorem}{Theorem}
\newtheorem{lemma}[theorem]{Lemma}
\newtheorem{remark}[theorem]{Remark}
\newcommand{\C}{\mathbb C}
\newcommand{\T}{\mathbb T}
\newcommand{\D}{\mathbb D}
\newcommand{\spoint}{\sigma_{\mathrm p}}
\DeclareMathOperator{\ran}{ran}
\title{The mixed point spectrum of a matrix contraction\\
on its closed unit ball}
\author{Alec Kriebel\thanks{Independent researcher.
ORCID: \href{https://orcid.org/0009-0001-9320-500X}{0009-0001-9320-500X}.}}
\date{October 5, 2026}

\begin{document}
\maketitle

\begin{abstract}
Let \(A\) be a contraction on a finite-dimensional complex normed space,
and let \(K_A f=f\circ A\) act on all continuous complex functions on its
closed unit ball. We show that every unimodular eigenfunction factors
through the contractive spectral projection onto the peripheral matrix
eigenspaces. Consequently the unimodular point spectrum is exactly the
multiplicative group generated by the peripheral matrix eigenvalues.
This answers the mixed point-spectrum converse expressly left open in
Kari Valentina K\"uster's 2015 thesis and included only in one direction
in her 2016 \emph{Pure Koopmanism} contribution. The argument is a short
finite-dimensional application of the classical Jacobs--de Leeuw--Glicksberg
mechanism. We give an elementary proof and collect the prior disk,
peripheral and nilpotent cases, together with their full-spectrum
consequences, in one table.
\end{abstract}

\section{The question and the classification}

Let \(X=(\C^k,\|\cdot\|)\), where \(k\geq1\) and \(\|\cdot\|\) is an
arbitrary complex norm. For a matrix \(A\) with induced operator norm
\(\|A\|\leq1\), put
\[
 U=\{x\in X:\|x\|\leq1\},\qquad
 K_A:C(U)\longrightarrow C(U),\qquad (K_Af)(x)=f(Ax).
\]
Here \(C(U)\) contains \emph{all} continuous complex-valued functions and
has the supremum norm. In particular \(U\) is compact, and \(\|K_A\|=1\)
because constants attain the contraction bound. Write
\[
 \T=\{\zeta\in\C:|\zeta|=1\},\qquad
 \D=\{\zeta\in\C:|\zeta|<1\}.
\]
The point spectrum \(\spoint(T)\) consists of those \(\zeta\) for which
\(Tf=\zeta f\) has a nonzero solution; \(\sigma(T)\) denotes the
Banach-operator spectrum.

K\"uster considered these exact matrix, norm, ball and observable
hypotheses in \cite[Section~3.1, printed pp.~35--43]{Kuester2015}
and \cite[Example~4, printed p.~321]{Kuester2016}. Her
Theorem~3.1.14(i), printed p.~43 (PDF p.~52), proves
\[
 \D\cup\langle\sigma(A)\cap\T\rangle
       \subseteq\spoint(K_A)
\]
when \(A\) has both a nonzero eigenvalue strictly inside the circle and a
peripheral eigenvalue. The paragraph following that theorem expressly
leaves the reverse inclusion open. Theorem~5(ii) of
\emph{Pure Koopmanism}, printed p.~321 (PDF p.~25), states the same
inclusion. The completion here is the exclusion of additional
unimodular eigenvalues.

To state all cases, define
\[
 S=\sigma(A)\cap\T,\qquad
 J=\{\alpha\in\sigma(A):0<|\alpha|<1\},
\]
\[
 \Gamma=\langle S\rangle
   =\left\{\prod_{\lambda\in S}\lambda^{m_\lambda}:
                      m_\lambda\in\mathbb Z\right\},
 \qquad \langle\varnothing\rangle=\{1\},
 \qquad
 Z_A=
 \begin{cases}
   \{0\},&A\text{ is singular},\\
   \varnothing,&A\text{ is invertible}.
 \end{cases}
\]
Negative exponents are essential in the definition of the group.
The bar on \(\Gamma\) below denotes closure in \(\C\).

\begin{theorem}\label{thm:classification}
Under the hypotheses above,
\[
\begin{array}{c|c|c}
\text{condition}&\spoint(K_A)&\sigma(K_A)\\ \hline
J\neq\varnothing&\D\cup\Gamma&\overline{\D}\\[2pt]
J=\varnothing&\Gamma\cup Z_A&\overline{\Gamma}\cup Z_A .
\end{array}
\]
In particular, in every case
\(\spoint(K_A)\cap\T=\Gamma\).
\end{theorem}

The all-peripheral point spectrum is prior work
\cite[Theorem~3.1.5, printed p.~38]{Kuester2015};
the nilpotent point spectrum is
\cite[Proposition~3.1.7, printed p.~39]{Kuester2015};
and the open disk from a nonzero stable matrix eigenvalue is
\cite[Theorem~3.1.13, printed pp.~41--42]{Kuester2015}.
The stable-only equality is stated in
\cite[Theorem~3.1.14(ii), printed p.~43]{Kuester2015} and
\cite[Theorem~5(iii), printed p.~321]{Kuester2016}.
The entire full-spectrum row for \(J\neq\varnothing\) is already an
immediate consequence of the prior open-disk inclusion, closedness of
spectrum, and \(\|K_A\|=1\).
The second full-spectrum row is a routine nilpotent direct-sum and
peripheral-rotation consequence, proved below for completeness.

\section{A projection compatible with the given norm}

\begin{lemma}\label{lem:projection}
There is an \(A\)-invariant decomposition \(X=E\oplus F\) and a commuting
projection \(P\) onto \(E\) such that
\[
 E=\bigoplus_{\lambda\in S}\ker(A-\lambda I),\qquad
 A|_E=B,\qquad A|_F=N,\qquad
 N^n\longrightarrow0,\qquad \|P\|\leq1.
\]
Moreover \(P(U)=V:=U\cap E\), and \(B\) is a surjective isometry on \(E\).
\end{lemma}

\begin{proof}
The bound \(\|A^n\|\leq1\) shows that \(\sigma(A)\) lies in the closed
disk and that every Jordan block associated with a unit-modulus
eigenvalue has size one: a larger block has powers with polynomial
growth. Take \(F\) to be the sum of the generalized eigenspaces associated
with eigenvalues of modulus less than one. The resulting spectral
projection \(P\) commutes with \(A\), and the Jordan formula gives
\(N^n\to0\) in operator norm.

We record the simultaneous recurrence needed to control \(P\) in the
\emph{original} norm. List the finitely many peripheral eigenvalues
as \(\lambda_1,\ldots,\lambda_r\). Compactness of \(\T^r\) supplies a
strictly increasing sequence \(m_j\) such that
\((\lambda_1^{m_j},\ldots,\lambda_r^{m_j})\) converges to a point of
\(\T^r\). For \(n_j=m_{2j}-m_j\), we have \(n_j\geq j\) and
\(\lambda_i^{n_j}\to1\) for every \(i\). If \(r=0\), simply take \(n_j=j\).
Since \(B\) is diagonalizable, this gives
\[
 B^{n_j}\longrightarrow I_E,\qquad A^{n_j}\longrightarrow P.
\]
Thus \(\|P\|\leq1\); this includes \(E=\{0\}\), where \(P=0\).
It follows that \(P(U)\subseteq V\), while the reverse inclusion follows
because \(P\) is the identity on \(E\).

When \(E\neq\{0\}\), \(B\) is invertible and
\(B^{n_j-1}\to B^{-1}\). Both \(B\) and \(B^{-1}\) have norm at most one
in the restricted norm on \(E\). Hence
\(\|Bx\|=\|x\|\), and \(B(V)=V\).
The zero-dimensional assertion has the same meaning on the singleton
\(V=\{0\}\).
\end{proof}

\begin{lemma}\label{lem:factor}
If \(f\in C(U)\) and \(K_Af=\zeta f\) with \(|\zeta|=1\), then
\[
 f=f\circ P.
\]
Pullback \(g\mapsto g\circ P\) identifies the unimodular eigenfunctions
of \(K_B\) on \(C(V)\) with those of \(K_A\) on \(C(U)\).
\end{lemma}

\begin{proof}
For \(x\in U\), both \(x\) and \(Px\) are in \(U\). Commutation of \(A\)
and \(P\), and the eigenfunction equation, give for every \(n\geq0\)
\[
 |f(x)-f(Px)|=|f(A^nx)-f(A^nPx)|.
\]
The difference of the two arguments is
\(N^n(I-P)x\), which tends to zero uniformly for \(x\in U\).
Uniform continuity of \(f\) on the compact ball therefore makes the
right side tend uniformly to zero. Thus \(f=f\circ P\).
Its restriction \(g=f|_V\) satisfies \(g(Bx)=\zeta g(x)\).
Conversely that equation pulls back because \(PA=BP\), and the
pullback of a nonzero \(g\) is nonzero since \(P(U)=V\).
\end{proof}

\section{The peripheral ball}

We give the familiar all-peripheral calculation directly on \(C(V)\).
If \(E=\{0\}\), the space \(V\) is a singleton and \(K_B\) is the identity
on constants, so both its spectra are \(\{1\}=\Gamma\).
Suppose henceforth that \(r=\dim E>0\). Choose complex-linear coordinates
\(z_1,\ldots,z_r\) diagonalizing \(B\), so
\[
 z_i(Bx)=\lambda_i z_i(x),\qquad |\lambda_i|=1.
\]
Every eigenvalue of \(K_B\) has modulus one, because
\(\|K_Bf\|_\infty=\|f\|_\infty\).
For nonnegative integer multi-indices \(p,q\), the monomial
\(z^p\overline z^{\,q}\) is an eigenfunction with eigenvalue
\[
 \gamma=\prod_{i=1}^r\lambda_i^{p_i-q_i}.
\]
It is nonzero on \(V\), since \(V\) contains a neighborhood of zero in
\(E\). These eigenvalues exhaust \(\Gamma\), proving
\(\Gamma\subseteq\spoint(K_B)\).

The polynomials in the coordinates and their conjugates are uniformly
dense in \(C(V)\) by the complex Stone--Weierstrass theorem: this algebra
contains constants, separates points and is closed under conjugation.
Let \(\zeta\in\T\setminus\Gamma\), and suppose \(K_Bf=\zeta f\).
Define
\[
 M_m=\frac1m\sum_{j=0}^{m-1}\zeta^{-j}K_B^j .
\]
Since \(B(V)=V\), \(K_B\) is an isometry and \(\|M_m\|\leq1\).
On each monomial eigenfunction the multiplier is
\(m^{-1}\sum_{j=0}^{m-1}(\gamma/\zeta)^j\), which tends to zero.
Consequently \(M_m h\to0\) for every such polynomial \(h\).
Given \(\varepsilon>0\), choose \(h\) with \(\|f-h\|_\infty<\varepsilon\).
As \(M_m f=f\), we have
\[
 \|f\|_\infty\leq\varepsilon+\|M_m h\|_\infty,
\]
and hence \(f=0\). No extra unimodular eigenvalue is possible.
Together with Lemma~\ref{lem:factor} this proves
\[
 \spoint(K_A)\cap\T=\spoint(K_B)=\Gamma.
 \tag{1}\label{eq:peripheral}
\]

For later use we also have
\[
 \sigma(K_B)=\overline\Gamma.
 \tag{2}\label{eq:peripheral-full}
\]
Indeed \(K_B\) and its inverse both have norm one, so their Neumann-series
resolvents exclude \(|\zeta|>1\) and \(|\zeta|<1\), respectively.
Thus \(\sigma(K_B)\subseteq\T\), while closedness gives
\(\overline\Gamma\subseteq\sigma(K_B)\).
If \(\Gamma\) is infinite, its closure is \(\T\). One elementary
justification is that an infinite subgroup has nonidentity elements with
arbitrarily small angular distance from \(1\); the successive powers of
each such element form an arbitrarily fine net of the circle.
If \(\Gamma\) is finite of order \(q\), it is the group of \(q\)-th roots
of unity and \(B^q=I\). Hence \(K_B^q=I\), and the factorization of
\(t^q-1\) excludes every other spectral value. This proves (2).

\section{Stable and nilpotent directions}

\subsection{A nonzero stable matrix eigenvalue}

Suppose \(\alpha\in J\), set \(a=|\alpha|\in(0,1)\), and choose a nonzero
complex-linear functional \(\ell\) with
\(\ell(Ax)=\alpha\ell(x)\). Such a left eigenfunctional exists for every
matrix eigenvalue, including when \(A\) has nontrivial Jordan blocks.
Let \(R=\max_{x\in U}|\ell(x)|>0\).
For \(0<|\mu|<1\), choose \(\theta\in\mathbb R\) with
\(e^{i\theta}=\mu/|\mu|\), and put
\[
 s=\frac{\log|\mu|+i\theta}{\log a},
 \qquad
 f_\mu(x)=
 \begin{cases}
 \exp\!\bigl(s\log(|\ell(x)|/R)\bigr),&\ell(x)\neq0,\\
 0,&\ell(x)=0.
 \end{cases}
\]
Here the logarithm in the function is the real logarithm of a positive
number. Since \(\Re s>0\), the function is continuous across
\(\ker\ell\), nonzero, and
\(f_\mu(Ax)=a^s f_\mu(x)=\mu f_\mu(x)\).
For \(\mu=0\), the nonzero continuous function
\[
 f_0(x)=\max\{0,|\ell(x)|/R-a\}
\]
vanishes on \(A(U)\), so \(K_Af_0=0\). This reconstructs the prior
open-disk inclusion by explicit radial eigenfunctions.
The norm bound excludes all eigenvalues outside \(\overline\D\);
(1) excludes the extra boundary values.
Therefore \(\spoint(K_A)=\D\cup\Gamma\).
Since spectrum is closed and contains the open disk,
\(\sigma(K_A)=\overline\D\), as already follows from the 2015 inclusion.

\subsection{Only peripheral eigenvalues and zero}

Suppose now that \(J=\varnothing\). All matrix eigenvalues on \(F\) are
zero, so \(N^d=0\) for some \(d\geq1\); if \(F=\{0\}\), choose \(d=1\).
The bounded projection
\[
 Q:C(U)\longrightarrow C(U),\qquad Qf=f\circ P
\]
commutes with \(K_A\). Its range is isometrically identified with \(C(V)\)
by pullback, and the induced restriction of \(K_A\) is \(K_B\).
Its kernel is exactly the functions vanishing on \(V\).
For \(f\in\ker Q\), the inclusion \(A^d(U)\subseteq V\) gives
\(K_A^df=0\). Thus
\[
 C(U)=\ran Q\oplus\ker Q
\]
is a closed invariant direct sum, with a peripheral restriction and a
nilpotent restriction.

When \(F\neq\{0\}\), a nonzero linear functional vanishing on \(E\)
restricts to a nonzero member of \(\ker Q\).
A nilpotent operator \(T\) on a nonzero space has zero as an eigenvalue:
choose \(h\neq0\) and the smallest positive \(m\) with \(T^m h=0\), so
\(T^{m-1}h\neq0\) lies in its kernel. Its spectrum is \(\{0\}\), since
for \(\zeta\neq0\) the finite geometric series in \(T/\zeta\) gives a
resolvent. When \(F=\{0\}\), the nilpotent summand is absent.
Spectra of a bounded operator on two closed invariant summands are the
unions of the restricted spectra, and the same holds for point spectra.
Combining (1)--(2) proves the second row of
Theorem~\ref{thm:classification}, because \(F\neq\{0\}\) here is equivalent
to \(A\) being singular. This completes the proof.

In particular every nilpotent matrix, including \(A=0\), gives
\(\spoint(K_A)=\sigma(K_A)=\{0,1\}\).
The optional zero in the second row requires \emph{singularity}, not that
\(A\) be the zero matrix.
In the first row zero is always an eigenvalue even if \(A\) is invertible:
for example \(A=\tfrac12 I\) maps \(U\) into a proper subset, and the
function \(f_0\) above lies in the Koopman kernel.

\section{Classical mechanism and scope of the contribution}

Lemma~\ref{lem:factor} is the elementary matrix-ball form of the classical
Jacobs--de Leeuw--Glicksberg reversible projection, rather than a new
general decomposition theorem. The connection can be checked explicitly.
Let \(H\) be the matrix-norm closure of \(\{A^n:n\geq0\}\).
It is a compact abelian semigroup. For every \(f\in C(U)\),
uniform continuity makes \(L\mapsto f\circ L\) continuous from \(H\)
to \(C(U)\). Thus the Koopman powers have compact strong operator closure,
which is also their weak operator closure.
Their composition order agrees with multiplication because \(H\) is
abelian.

The closure \(G_B\) of the powers of the surjective isometry \(B\) is a
compact group. Recurrence and \(N^n\to0\) show that
\[
 G=\{DP:D\in G_B\}\subseteq H
\]
is a group ideal with identity \(P\). Every nonempty ideal of \(H\)
intersects \(G\), by multiplication by \(P\); its intersection with this
group is then the whole group. Hence \(G\) is the minimal ideal, and \(P\)
is its identity idempotent. The corresponding operator idempotent is
exactly \(Qf=f\circ P\).
The classical JdLG theorem places every unimodular eigenvector in
\(\ran Q\), giving the factorization; see
\cite[construction on printed pp.~308--309 and
Theorem~16.33(a), printed p.~313]{EFHN}.
Those references were checked in the author manuscript dated
April~22, 2016, of the book published in 2015.
The norm compatibility of \(P\) and the identification of \(Q\) above
are the concrete applicability steps in this setting.

The theorem therefore answers a question expressly left open in the
2015 source and stated only as an inclusion in the 2016 source, by a
short specialization of an existing mechanism. A bounded source audit
located no earlier explicit completion in its examined scope; this is
not a claim of historical firstness, exhaustive literature coverage, or
continuous openness from 2015 to 2026. The accompanying priority note
records the read scope and remaining coverage gaps.

All assertions concern the specified closed ball and the full continuous
observable space. No invariant measure, holomorphic restriction,
smoothness assumption, matrix diagonalizability assumption, or matrix
invertibility assumption enters the theorem. Finite algebraic checks
can test examples, projections and eigenfunction identities, but do not
prove the infinite-dimensional spectral classification; the proof is the
argument above.

\paragraph{AI assistance and review status.}
Extensive AI assistance was used in solving, drafting, literature
comparison, and adversarial verification of this work. Separate AI
verification tasks examined the peripheral projection, the spectral
cases, and finite computational controls. These checks are not
conventional human peer review. This research note has not undergone
conventional human peer review.

\begin{thebibliography}{9}

\bibitem{Kuester2015}
Kari Valentina K\"uster,
\emph{The Koopman Linearization of Dynamical Systems},
Diplomarbeit, Eberhard Karls Universit\"at T\"ubingen, March 2015.
Relevant scope: Section~3.1, printed pp.~35--43 (PDF pp.~44--52);
Theorem~3.1.14(i) and its open-converse paragraph, printed p.~43
(PDF p.~52).
\url{https://homepages.laas.fr/henrion/mfo16/kari-kuester.pdf}.

\bibitem{Kuester2016}
Kari K\"uster,
\emph{Pure Koopmanism},
in \emph{Mini-Workshop: Applied Koopmanism},
Oberwolfach Report 7/2016, printed pp.~320--322
(PDF pp.~24--26).
Relevant statement: Example~4 and Theorem~5, printed p.~321
(PDF p.~25).
\url{https://ems.press/content/serial-article-files/46612}.

\bibitem{EFHN}
Tanja Eisner, B\'alint Farkas, Markus Haase, and Rainer Nagel,
\emph{Operator Theoretic Aspects of Ergodic Theory},
Graduate Texts in Mathematics 272, Springer, 2015.
The checked source is the author manuscript dated April~22, 2016:
JdLG construction, printed pp.~308--309 (PDF pp.~326--327);
Theorem~16.33(a), printed p.~313 (PDF p.~331).
\url{https://www.math.uni-leipzig.de/~eisner/book-EFHN.pdf}.

\end{thebibliography}
\end{document}
